\documentclass[aspectratio=169]{beamer}
\usepackage{udec-slides}
\ayud{2}{Campo Eléctrico}

\begin{document}
\primeraDiapo

\begin{frame}{Campo Eléctrico}{Expresión para una Carga Puntual}
    \begin{boxTeoria}{Campo Eléctrico debido a Carga Puntual}
        Si tenemos una carga fuente $Q$, el campo que genera sobre un punto
        $P$ es:
        \begin{equation*}
            \vec{E} = \frac{1}{4\pi\epsilon_0} \frac{Q}{|\vec{R}|^2} \hat{R}
        \end{equation*}

        Donde:
        \begin{itemize}
            \item $\vec{R} = \vec{r}_P - \vec{r}_Q$
        \end{itemize}
    \end{boxTeoria}
\end{frame}

\begin{frame}{Problema 1}{Distribución de Carga Lineal}

    Considere un hilo delgado en forma de semicircunferencia de radio $R$ cargado con una densidad de carga lineal uniforme $\lambda$. Este sitúa su centro el origen. Determine, el vector campo eléctrico $\vec{E}$ generado por esta distribución en el punto $P$ ubicado en el origen.

    \begin{center}
        \begin{tikzpicture}[scale=2, >=stealth, y=1cm, x=1cm]
            \draw[->, thin, gray] (-1.3, 0) -- (1.3, 0) node[right, black, scale=0.7] {$x$};
            \draw[->, thin, gray] (0, -0.2) -- (0, 1.3) node[above, black, scale=0.7] {$y$};

            \draw[ultra thick, UdeCPrincipal] (1,0) arc (0:180:1);

            \filldraw[black] (0,0) circle (1.2pt) node[below left, scale=0.8] {$P$};

            \draw[dashed, red, thick] (0,0) -- (0.707, 0.707) node[midway, above left, scale=0.8, red] {$R$};

            \node[UdeCPrincipal, scale=0.9] at (135:1.3) {$\lambda$};

        \end{tikzpicture}
    \end{center}

\end{frame}

\begin{frame}{Problema 2}{Distribución de Carga Superficial}

    Considere un plano infinito que coincide con el plano $xy$ ($z = 0$), el cual posee una densidad de carga superficial uniforme $\sigma$. Exprese la integral que permite calcular la magnitud del campo eléctrico que esta ejerce sobre una carga ubicada en $P=h\hat{z}$.


\end{frame}

\begin{frame}{Problema 3}{Distribución de Carga Volumétrica}
    Una esfera sólida de radio $a$ posee una densidad de carga volumétrica no uniforme definida por:

    \begin{equation*}
        \rho(r) = \rho_0 \left( \frac{r - a}{a} \right)
    \end{equation*}

    donde $r$ es la distancia radial desde el centro de la esfera. Encontrar la intensidad del campo eléctrico en un punto situado a una distancia $R$ muy grande ($R \gg a$). ¿Qué cosas deben asumirse?
\end{frame}

\end{document}